\documentclass[11pt,a4paper]{article}
\usepackage[utf8]{inputenc}
\usepackage{amsmath,amssymb,amsfonts}
\usepackage{geometry}
\usepackage{hyperref}
\usepackage{xcolor}

\geometry{margin=1in}

\title{\textbf{\Huge Hermite 4-Point Cubic Interpolation \& Varispeed Resampling}\\
\Large DSP Mathematics of Wavetable \& Sampler Objects}
\author{\textbf{Time Dilation DAW (v4.0) Engineering Group}}
\date{\today}

\begin{document}

\maketitle

\begin{abstract}
This paper formulates the mathematical models of sub-sample fractional position interpolation used in \textbf{Time Dilation DAW}'s \texttt{[osc\textasciitilde{}]} and \texttt{[sampler\textasciitilde{}]} nodes, including 4-Point Hermite Cubic Splines, Linear Interpolation, and Nearest-Neighbor Zero-Order Hold.
\end{abstract}

\section{Fractional Table Read Head Position}
When playing audio buffers or wavetables at non-integer pitch ratios or dilated time speeds $\gamma$, the fractional read head position at sample $n$ is:
\begin{equation}
p(n) = p(n-1) + \text{step}, \quad \text{where } \text{step} = \frac{|\gamma| \cdot f \cdot N_{\text{table}}}{f_s}
\end{equation}
Decomposing $p(n)$ into integer index $i = \lfloor p(n) \rfloor$ and fractional offset $\alpha = p(n) - i \in [0, 1)$.

\section{Hermite 4-Point Cubic Spline Interpolation}
To eliminate high-frequency aliasing and inter-sample quantization noise during extreme varispeed pitch shifting, we evaluate four adjacent table samples $(y_0, y_1, y_2, y_3)$ at indices $(i-1, i, i+1, i+2)$:
\begin{equation}
y(\alpha) = c_3 \alpha^3 + c_2 \alpha^2 + c_1 \alpha + c_0
\end{equation}
where the polynomial coefficients are computed as:
\begin{align}
c_0 &= y_1 \\
c_1 &= \frac{1}{2}(y_2 - y_0) \\
c_2 &= y_0 - \frac{5}{2} y_1 + 2 y_2 - \frac{1}{2} y_3 \\
c_3 &= \frac{1}{2}(y_3 - y_0) + \frac{3}{2}(y_1 - y_2)
\end{align}

\section{Linear Interpolation \& Nearest Neighbor}
\subsection{Linear Interpolation (2-Point)}
\begin{equation}
y(\alpha) = (1 - \alpha) y_1 + \alpha y_2
\end{equation}
\subsection{Nearest Neighbor (Zero-Order Hold)}
\begin{equation}
y(\alpha) = y_{\lfloor \alpha + 0.5 \rfloor + i}
\end{equation}

\end{document}
